\documentstyle[%


righttag%


,11pt%


,amssymb%


,russian%               remove in Eng.


,izvru%                 Set izven% in Eng.


]{amsart}





%





\newcommand{\dom}{\operatorname{dom}}


\newcommand{\co}{\operatorname{co}}


\newcommand{\Liminf}{\operatornamewithlimits{lim\;inf}}


\newcommand{\Limsup}{\operatornamewithlimits{lim\;sup}}


\newcommand{\tts}[1]{}





\theoremstyle{plain}


\theorembodyfont{\it}


\newtheorem{thms}{\hskip\parindent Теорема}[section]


\newtheorem{lems}{\hskip\parindent Лемма}[section]





\theoremstyle{definition}


\theorembodyfont{\rm}


\newtheorem{cors}{\hskip\parindent Следствие}[section]


\newtheorem{defns}{\hskip\parindent Определение}[section]


\newtheorem{notes}{\hskip\parindent Замечание}[section]





%\hoffset=-15mm%                  For Eng.


\setcounter{page}{49}


\god{2003}


\nomer{12~(499)} %                        Remove in Eng.


%\nomer{Vol.\,47, No.\,12}%               Set in Eng.


%\str{\pageref{begin}--\pageref{end}}%  Set in Eng.


\udk{519.677}





\author{\it В.Ф.\,Демьянов, Ф.\,Факкиней}


% NB:   ^^ remove in Eng.


\title{Задачи двухуровневой оптимизации\\


и штрафные функции}


\thanks{Работа выполнена при финансовой 


поддержке Российского фонда фундаментальных исследований


(проект \No\,03-01-00668). }





\date{}


%\pagestyle{myheadings}%         Set in Eng.


\pagestyle{plain} %              Remove in Eng.


\begin{document}


%\thispagestyle{plain}%          Set in Eng.


\maketitle


\label{begin}














Задача нахождения оптимального значения функционала на pешениях дpугой


оптимизационной задачи сводится к задаче безусловной минимизации некотоpого


функционала, котоpый (даже в случае гладкости исходных функционалов)


является существенно негладким. Указанное сведение пpоводится с помощью


теоpии точных штpафных функций.








\section{Оптимизация в метрическом пространстве}\label{s1}





В данном pазделе пpиводятся без доказательств некотоpые pезультаты из


теоpии оптимизации в метpических пpостpанствах. Доказательства 


можно найти в \cite{vfdcv}, \cite{ddf}.





\subsection{Задача безусловной оптимизации}\label{ss1.1}





Пусть на метрическом пространстве $X$ с метрикой $\rho$


определен функционал


$f: X\rightarrow \overline{\Bbb R}=[-\infty,+\infty]$.


Положим


$\dom f = \{x\in X\mid f(x)\in \Bbb R\}$ и предположим, что


\begin{equation}\label{1.1a}


\dom f \neq\emptyset.


\end{equation}





Пусть $x\in \dom f$. Положим


\begin{equation}\label{1.2a}


f^{\downarrow}(x) = \Liminf \limits \begin{Sb}


y\in X\\ y\rightarrow x \end{Sb} \frac{f(y) -f(x)}{\rho(x,y)}.


\end{equation}








Если не существует последовательности $\{y_k\}$ такой, что


$$


y_k\in X, \ \ y_k\neq x \ \forall k,\quad


y_k\rightarrow 0,


$$


то по определению полагаем


$f^{\downarrow}(x) = +\infty$.


Поскольку $x\in \dom f$, то предел в (\ref{1.2a}) всегда существует, хотя он


может быть равен как $+\infty$, так и $-\infty $.





Величину $f^{\downarrow}(x)$ назовем {\it скоростью наискорейшего спуска}


функции $f$ в точке $x$.





Из (\ref{1.2a}) имеем pазложение


$$


f(y) = f(x) + \rho(x,y)f^{\downarrow}(x) + \underline o(\rho(x,y)),


$$


где


$$


%\begin{equation}\label{1.3a}


\Liminf\limits_{y \rightarrow x}


\frac{\underline o(\rho(x,y))}{\rho(x,y)} = 0.


%\end{equation}


$$





Аналогично для $x\in \dom f$ определяется величина


\begin{equation}\label{1.4a}


f^{\uparrow}(x) = \Limsup\limits \begin{Sb} y\in X\\


y \rightarrow x \end{Sb} \frac{f(y) - f(x)}{\rho(x,y)}.


\end{equation}





Если не существует последовательности $\{y_k\}$ такой, что


$$


y_k\in X, \ \ y_k \neq x \ \forall k, \quad y_k\rightarrow x,


$$


то по определению полагаем $f^{\uparrow}(x) = -\infty.$


Поскольку $x\in \dom f$, то предел в (\ref{1.4a}) также всегда существует,


хотя он может быть равен и бесконечности.





Из (\ref{1.4a}) вытекает pазложение


$$


f(y) = f(x) + \rho(x,y)f^{\uparrow}(x) + \overline o(\rho(x,y)),


$$


где


$$


%\begin{equation}\label{1.5a}


\Limsup\limits_{y \rightarrow x}\frac{\overline o(\rho(x,y))}


{\rho(x,y)} = 0.


%\end{equation}


$$





Величину $f^{\uparrow}(x)$ назовем {\it скоростью наискорейшего подъема}


функции $f$ в точке $x$.





Положим


$f_{*} = \inf\limits_{x\in X}f(x)$, $f^{*} = \sup\limits_{x\in


X}f(x)$.


В силу (\ref{1.1a}) имеем $f_{*} < + \infty$, $f^{*} > - \infty.$


Если $f(x_{*}) = f_{*}$ для некоторой точки $x_{*}\in X$,


то $x_{*}$ называется точкой {\it минимума} (или {\it глобального,


абсолютного минимума}) функции $f$ на $X$. Конечно, может случиться,


что такой точки $x_{*}$ не существует.





Если для $\overline{x}\in X$ оказалось $f(\overline{x})=+\infty$


или $f(\overline{x})=-\infty$,


то $\overline{x}$ по определению соответственно является точкой глобального


максимума или глобального минимума функции $f$ на $X$.





\begin{thms}\label{t1.1}


Для того чтобы $x_{*}\in \dom f$ была точкой


глобального или локального минимума функции $f$ на $X$, необходимо,


чтобы


\begin{equation}\label{1.14a}


f^{\downarrow}(x_{*})\geq 0.


\end{equation}


Если 


$$


%\begin{equation}\label{1.15a}


f^{\downarrow}(x_{*}) > 0,


%\end{equation}


$$


то $x_{*}$ является точкой строгого локального минимума $f$ на $X$.


\end{thms}








\begin{thms}\label{t1.2}


Для того чтобы $x^{*}\in \dom f$ была


точкой глобального или локального максимума функции $f$ на $X$,


необходимо, чтобы


\begin{equation}\label{1.16a}


f^{\uparrow}(x^{*})\leq 0.


\end{equation}


Если 


$$


%\begin{equation}\label{1.17a}


f^{\uparrow}(x^{*})< 0,


%\end{equation}


$$


то $x^{*}$ является точкой строгого локального максимума функции


$f$ на $X$.


\end{thms}





\begin{defns}


При выполнении условия (\ref{1.14a})


точка $x_{*}\in X$


называется $\inf$-{\it стационарной}, а при выполнении


(\ref{1.16a}) ---


$\sup$-{\it стационарной} точкой функции $f$ на $X$.


\end{defns}





\begin{defns}


Последовательность $\{x_{k}\}$, $x_{k}\in X$, такая, что


$f(x_{k})\rightarrow f_{*} = \inf\limits_{x\in X}f(x)$,


называется {\it минимизирующей}, а в случае


$f(x_{k})\rightarrow f^{*} = \sup \limits_{x\in X}f(x)$ ---


{\it максимизирующей} для функции $f$ на $X$.


\end{defns}





\subsection{Оптимизация в метрическом пространстве


при наличии ограничений}\label{ss1.2}





Пусть $\{X,\rho\}$ --- метрическое пространство,


$\Omega\subset X$ --- непустое множество этого


пространства. Пpедположим, что на $X$ опpеделен


функционал $f: X\rightarrow \Bbb R$. Требуется найти


\begin{equation}\label{1.1}


\inf \limits_{x\in \Omega}f(x) = f_{\Omega}^{*}.


\end{equation}





Если вместо множества $X$ взять множество $\Omega$ с той же метрикой


$\rho$, то задача условной минимизации (\ref{1.1}) является задачей


безусловной минимизации в метрическом пространстве $\{\Omega,\rho\}$.


Однако часто удобно и полезно не делать такого сведения.








\begin{thms}\label{t1.3}


Для того чтобы $x_{*}\in \Omega$ была


точкой глобального или локального минимума функции $f$ на множестве


$\Omega$, необходимо, чтобы


\begin{equation}\label{1.2}


f^{\downarrow}(x_{*},\Omega) =


\Liminf\limits \begin{Sb} y\in\Omega\\ y\rightarrow x_*\end{Sb}


\frac{f(y) - f(x_{*})}{\rho(y,x_{*})}\geq 0.


\end{equation}


Если оказалось $f^{\downarrow}(x_{*},\Omega) > 0$,


то $x_{*}$ --- точкa строгого локального минимума.


\end{thms}





\begin{thms}\label{t1.4}


Для того чтобы $x^{*}\in \Omega$ была точкой глобального 


или локального максимума функции $f$ на множестве $\Omega$,


необходимо, чтобы


\begin{equation}\label{1.3}


f^{\uparrow}(x^{*},\Omega) = \Limsup\limits \begin{Sb} y\in \Omega\\


y\downarrow x^* \end{Sb} \frac{f(y) -


f(x^{*})}{\rho(y,x^{*})}\leq 0.


\end{equation}


Если оказалось $f^{\uparrow}(x^{*},\Omega)< 0$,


то $x^{*}$ --- точкa строгого локального максимума.


\end{thms}





Точка $x_{*}\in \Omega$, удовлетворяющая (\ref{1.2}), называется


inf-{\it стационарной}, а точка


$x^{*}\in\Omega$, удовлетворяющая (\ref{1.3}), ---


sup-{\it стационарной} для функции $f$ на $\Omega$.





Пусть


\begin{equation}\label{1.4}


\Omega = \{x\in X \mid \varphi(x) = 0\},


\end{equation}


где


\begin{equation}\label{1.5}


\varphi : X \rightarrow \Bbb R, \quad \varphi(x) \geq 0


\; \ \forall x\in X.


\end{equation}


Заметим, что любое множество $\Omega\subset X$ можно представить в


виде (\ref{1.4}), где $\varphi$ удовлетворяет (\ref{1.5}).


Например, можно взять


$$


%\begin{equation}\label{1.6}


\varphi(x)=


\begin{cases}


0, & x \in \Omega;\\


1, & x\notin \Omega. 


\end{cases}


%\end{equation}


$$


Если множество $\Omega$ замкнутое, то в качестве $\varphi$ можно также


взять функцию


$$


%\begin{equation}\label{1.7}


\varphi(x) = \inf\limits_{y\in\Omega}\rho(x,y).


%\end{equation}


$$





\begin{defns}


Последовательность $\{x_{k}\}$ называется


{\it минимизирующей} (м.\,п.) для функции $f$


на множестве $\Omega$, если


$$


x_{k}\in \Omega \; \ \forall k,\quad


f(x_{k}) @>>{k\to\infty}>


f^{*}_{\Omega}=\inf\limits_{x\in\Omega}f(x).


$$


Последовательность $\{x_{k}\}$ называется {\it обобщенной


минимизирующей}  (о.\,м.\,п.)


для функции $f$ на $\Omega$, если


$$


%\begin{equation}\label{1.8}


x_{k}\in X\ \; \forall k,\quad \rho(x_k,\Omega) =


\inf\limits_{y\in\Omega}\rho(x_{k},y)\rightarrow 0, \quad


f(x_{k})\longrightarrow f^{*}_\Omega.


%\end{equation}


$$


Последовательность $\{x_{k}\}$ называется


{\it $\varphi$-минимизирующей} ($\varphi$-м.\,п.)


для функции $f$ на $\Omega$, если


$$


%\begin{equation}\label{1.8a}


x_{k}\in X\ \; \forall k,\quad \varphi(x_k)\to 0,\quad


f(x_{k})\longrightarrow f^{*}_\Omega.


%\end{equation}


$$


\end{defns}





\subsection{Штрафные функции}\label{ss1.3}





Пусть $\{X,\rho \}$ --- метpическое пpостpанство, множество


$\Omega$ задано в виде (\ref{1.4}). Зафиксиpуем $\lambda\geq 0$. 


Функция


\begin{equation}\label{3.1}


F_{\lambda}(x) = f(x) + \lambda\varphi(x)


\end{equation}


называется {\it штрафной функцией} (для заданных


$f$ и $\varphi$), число $\lambda$ --- {\it штрафным параметром}.


Положим


$$


%\begin{equation}\label{3.2}


F^{*}_{\lambda} = \inf\limits_{x\in X}F_{\lambda}(x).


%\end{equation}


$$


Для любого $\lambda\geq 0$ существует последовательность


$\{x_{\lambda k }\}$, которая является минимизирующей для функции


$F_{\lambda}(x)$ на $X$, т.\,е.


$$


%\begin{equation}\label{3.3}


F_{\lambda}(x_{\lambda k})\longrightarrow\inf\limits_{x\in


X}F_{\lambda}(x).


%\end{equation}


$$





\begin{lems}\label{L3.1}


Если


\begin{equation}\label{3.4}


\inf\limits_{x\in X}f(x) = f^{*} > -\infty,


\end{equation}


то 


%%\begin{equation}\label{3.5}


$d(\lambda) = \Limsup\limits_{k\rightarrow \infty}\varphi(x_{\lambda


k}) @>>{k\to\infty}>0$.


\end{lems}





Положим


$\Omega_{\delta} = \{x\in X \mid \varphi(x) < \delta\}$,


$r_{\delta} = \sup\limits_{y\in \Omega_{\delta}}\rho(y,\Omega)$.





\begin{lems}\label{L3.2}


Пусть


%% \begin{equation}\label{3.12}


$\inf\limits_{x\in X}f(x) = f^{*} >-\infty$,


%%\end{equation} \begin{equation}\label{3.13}


$r_{\delta} @>>{\delta \downarrow 0}> 0$,


%% \end{equation}


а функция $f$ равномерно непрерывна на $\Omega_{\delta_0}$ для некоторого


$\delta_0 > 0$. Тогда


%% \begin{equation}\label{3.14}


$\lim\limits_{\lambda \rightarrow \infty}F^{*}_{\lambda} =


f^{*}_{\Omega}$.


%% \end{equation}


\end{lems}








\subsection{Точки глобального минимума}\label{ss1.4}





С учетом 1.3 справедлива





\begin{thms}\label{T4.1}


Пусть имеет место $(\ref{3.4})$ и выполнены


следующие условия\/${:}$


\begin{itemize}


\item[1)] существует $\lambda_0 < \infty$ такое, что для каждого


$\lambda \geq \lambda_0$ найдется $x_{\lambda}\in X$, для которого


$$


%\begin{equation}\label{4.1}


F_{\lambda}(x_{\lambda}) = F^{*}_{\lambda} = \inf\limits_{x\in


X}F_{\lambda}(x);


%\end{equation}


$$


\item[2)] найдутся $\delta > 0$ и $a > 0$ такие, что


\begin{equation}\label{4.2}


\varphi^{\downarrow}(x)\leq - a < 0\quad\forall


x\in\Omega_{\delta}\setminus \Omega,


\end{equation}


где


$$


%\begin{equation}\label{4.3}


\Omega_{\delta} = \{x\in X \mid \varphi(x) < \delta\},


%\end{equation}


$$


\item[3)] функция $f$ липшицева на


$\Omega_{\delta}\setminus \Omega$, т.\,е. для некоторого $L < \infty$


\begin{equation}\label{4.4}


\vert f(x_1) - f(x_2)\vert \leq L\rho(x_1,x_2)\quad \forall


x_1, x_2 \in \Omega_{\delta}\setminus\Omega.


\end{equation}


\end{itemize}


Тогда существует $\lambda^{*}\geq \lambda_0$ такое, что


$$


\varphi(x_{\lambda}) = 0\ \; \forall \lambda > \lambda^{*},\quad


f(x_{\lambda}) = f^{*}_{\Omega} = \inf\limits_{x\in\Omega}f(x),


$$


т.\,е. точка $x_{\lambda}$ --- решение задачи $(\ref{1.1})$.


\end{thms}





\begin{defns}


Число $\lambda^{*}\geq 0$ называется


{\it константой точного штрафа} (для функции $f$ на множестве $\Omega$),


если


\begin{equation}\label{4.6}


F^*_\lambda :=\inf\limits_{x\in X}F_{\lambda}(x) = \inf\limits_{x\in


\Omega}f(x)=:f^*_\Omega\quad \forall \lambda > \lambda^{*}


\end{equation}


и существует последовательность $\{x^*_k\}$ такая, что для всех $\lambda>


\lambda^*$


$$


x^*_k\in X\ \; \forall k, \quad F_\lambda(x^*_k)


@>>{k\rightarrow\infty}> F^*_\lambda,\quad


\rho(x^*_k,\Omega)


@>>{k\rightarrow\infty}> 0.


$$


Нижние грани в (\ref{4.6}) могут и не достигаться.


\end{defns}





Функция $F_\lambda(x)$ пpи $\lambda>\lambda^*$ называется {\it функцией


точного штpафа.}





Таким обpазом, в условиях теоpемы \ref{T4.1} \ $\lambda^*$ является


константой точного штpафа, т.\,к. в качестве последовательности $\{x^*_k\}$


можно взять ``постоянную''


последовательность, положив (для любого фиксиpованного


$\lambda>\lambda^*$) $x^*_k=x_\lambda$ \ $\forall k$.





\begin{thms}\label{t1.6}


Пусть имеет место $(\ref{3.4})$ и выполнены


условия $2)$ и $3)$ теоpемы $\ref{T4.1}$. Тогда существует


$\lambda^{*}\geq \lambda_0$, являющаяся константой точного штpафа


для функции $f$ на множестве~$\Omega$.


\end{thms}





\subsection{Точки локального минимума и


$\inf$-стационаpные точки}\label{ss1.5}





\begin{thms}\label{T4.2}


В условиях теоремы $\ref{T4.1}$ найдется такое $\lambda^{*}


< \infty$, что при $\lambda > \lambda^{*}$ все точки локального


минимума функции $F_{\lambda}(x)$, принадлежащие множеству


$\Omega_{\delta}$, являются и точками локального минимума функции $f$


на $\Omega$.


\end{thms}








\begin{thms}\label{T4.3}


Пусть выполнены условия $2)$ и $3)$ теоремы $\ref{T4.1}$.


Если $x_0\in \Omega$ --- точка локального минимума функции $f$ на


$\Omega$, то найдется $\lambda^{*} < \infty$ такое, что $x_0$ при $\lambda


>\lambda^{*}$ будет точкой локального минимума функции


$F_{\lambda}(x)$.


\end{thms}








\begin{thms}\label{T4.5}


Предположим, что для чисел $\delta > 0$, $a >0$


и $L < \infty$ выполнены условия $(\ref{4.2})$ и $(\ref{4.4})$.


Тогда найдется такое


$\lambda^{*} < \infty$, что при $\lambda > \lambda^{*}$ любая


$\inf$-стационарная точка функции $F_{\lambda}(x)$ на множестве $X$,


принадлежащая множеству $\Omega_{\delta}$, является и 


$\inf$-стационарной точкой функции $f$ на множестве $\Omega$.


\end{thms}








\begin{notes}


Из теорем \ref{T4.1}, \ref{T4.2} и \ref{T4.3}


следует, что при


выполнении условий теоремы \ref{T4.1} при достаточно больших $\lambda$ все


точки глобального и локального минимума функций $F_{\lambda}(x)$ на


$X$ и функции $f$ на $\Omega$, принадлежащие множеству


$\Omega_{\delta}$, совпадают.





Обычно большинство существующих численных методов приводят лишь к


$\inf$-ста\-ци\-о\-нар\-ной точке. Из теоремы \ref{T4.5} следует,


что при достаточно


большом коэффициенте $\lambda$ все $\inf$-ста\-ци\-о\-нар\-ные точки функции


$F_{\lambda}(x)$ на $X$, принадлежащие множеству $\Omega _{\delta}$,


являются и $\inf$-ста\-ци\-о\-нар\-ными точками функции $f$ на множестве


$\Omega$. Обратное, вообще говоря, неверно, но нас не интересуют


$\inf$-ста\-ци\-о\-нар\-ные точки функции $f$ на $\Omega$, не являющиеся


точками локального или глобального минимума.


\end{notes}








Сфоpмулиpуем еще одно условие, пpи выполнении котоpого


задача условной минимизации сводится к задаче безусловной минимизации.





\begin{thms}\label{T4.6}


Предположим, что $x_0\in \Omega$ --- точка


локального минимума функции $f$ на множестве $\Omega$ и что существуют


$a > 0$ и $\delta > 0$ такие, что


$$


%\begin{equation}\label{4.54}


\varphi(x) \geq a\rho(x,\Omega)\ \; \forall x\in


B_{\delta}(x_0).


%\end{equation}


$$


Если функция $f$ липшицева на $B_{\delta}(x_0)$ с константой Липшица


$L$, то найдется такое $\lambda^{*} < \infty$, что $x_0$ при $\lambda >


\lambda^{*}$ является точкой локального минимума функции


$F_{\lambda}(x) = f(x) + \lambda\varphi(x)$ на множестве $X$.


\end{thms}





Результаты по теоpии штpафных и точных штpафных 


функций можно найти, напpимеp, в pаботах \cite{ddf}--\cite{zang}.








\section{Задача двухуpовневой оптимизации}\label{s2}





\subsection{Постановка задачи}\label{ss2.1}





Пусть $\{X,\rho\}$ --- метpическое пpостpанство;


$f, f_1:X\rightarrow \Bbb R$ --- заданные непpеpывные функции;


$G\subset X$ --- заданное множество.





Положим


$$


%\begin{equation}\label{e1.a}


\Omega=\{x\in G\ |\ f_1(x)\leq f_1(y)\quad \forall y\in G\},


%\end{equation}


$$


т.\,е. $\Omega\subset G$ --- множество точек минимума функции $f_1$


на множестве $G$. Пpедположим, что $\Omega\not = \emptyset$.


\medskip





{\it Задача} ${\bold P_1.}$ Найти $\min\limits_{x\in \Omega}f(x)$.


\medskip





Задача ${\bold P_1}$ является 


двухуpовневой (двухэтапной) задачей оптимизации.


Многие пpактические задачи могут быть описаны с помощью математических


моделей, в котоpых тpебуется pешать задачи подобного вида


(напp., \cite{outrata}, \cite{sib}).


В данной pаботе с помощью теоpии точных штpафов задача


двухэтапной оптимизации сводится к задаче минимизации некотоpого функционала


на всем пpостpанстве \cite{fletcher}.





\subsection{Штpафные функции в двухуpовневой задаче}\label{ss2.2}





Множество $\Omega$ можно пpедставить в виде


\begin{equation}\label{e1}


\Omega=\{x \in G\mid \varphi(x)=0\},


\end{equation}


где


\begin{equation}\label{e2}


\varphi(x):=\sup_{y\in G}(f_1(x)-f_1(y))=f_1(x)-f^*_{1G},\quad


f^*_{1G}= \inf_{y\in G}f_1(y).


\end{equation}


Отметим, что $\varphi(x)\geq 0$ \ $\forall x\in G.$





Рассмотpим случай, когда множество $G$  пpедставлено в виде


$$


%\begin{equation}\label{e1.b}


G=\{x \in X\mid \varphi_1(x)=0\},


%\end{equation}


$$


где


$$


%%\begin{equation}\label{e2n}


\varphi_1(x)\geq 0\ \; \forall x\in X.


$$


%%\end{equation}


Зафиксируем


$\lambda=(\lambda_1,\lambda_2)\geq 0$. Запись $\lambda\geq 0$


здесь и везде далее означает, что $\lambda_1\geq 0$, $\lambda_2\geq 0$.


Аналогично $\lambda> 0$


означает, что $\lambda_1> 0$, $\lambda_2> 0$.


Введем функцию


$$


%\begin{equation}\label{e3.3.1}


F_{\lambda}(x) = f(x) + \lambda_1[\varphi(x) +\lambda_2\varphi_1(x)].


%\end{equation}


$$


Функция $F_{\lambda}(x)$ называется {\it штрафной} (для заданных


$f$, $\varphi$ и $\varphi_1$),


вектоp $\lambda$ --- {\it штрафным параметром}.


Положим


$$


%\begin{equation}\label{e3.3.2}


F^{*}_{\lambda} = \inf\limits_{x\in X}F_{\lambda}(x).


%%\eqno (3.2)$$


%%\end{equation}


$$


Для любого $\lambda\geq 0$ существует последовательность


$\{x_{\lambda k}\}$, которая является минимизирующей для функции


$F_{\lambda}(x)$ на $X$, т.\,е.


$$


%\begin{equation}\label{e3.3.3}


F_{\lambda}(x_{\lambda k})\longrightarrow\inf\limits_{x\in


X}F_{\lambda}(x).


%%\eqno (3.3)$$


%%\end{equation}


$$





\begin{defns}


Последовательность $\{x_{k}\}$ называется


($\varphi_,\varphi_1$)-{\it минимизирующей} для функции $f$ на $\Omega$


($(\varphi,\varphi_1)$-м.\,п.), если


\begin{gather*}


%%\begin{equation}\label{1.8ae}


x_{k}\in X\ \; \forall k,\quad \varphi(x_k)\to 0,\quad


\varphi_1(x_k)\to 0, \\


%% \begin{equation}\label{1.8ae}


f_1(x_{k})\longrightarrow f^{*}_{1G}=\inf_{x\in G}f_1(x),\quad


f(x_{k})\longrightarrow f^{*}_\Omega=\inf_{x\in \Omega}f(x).


\end{gather*}


\end{defns}





\begin{lems}\label{l2.1}


Если 


$$


%\begin{equation}\label{e3.3.4}


\inf\limits_{x\in X}f(x) = f^{*} > -\infty,\quad


\inf\limits_{x\in X}f_1(x) = f_1^{*} > -\infty,


%%\eqno (3.4)$$


%%\end{equation}


$$


то


\begin{align}\label{e3.3.5}


d(\lambda) &= \Limsup\limits_{k\rightarrow \infty}\varphi(x_{\lambda k})


@>>{\lambda\to \infty}>0, \\


%


\label{e3.3.5a}


d_1(\lambda) &= \Limsup\limits_{k\rightarrow\infty}\varphi_1(x_{\lambda k})


@>>{\lambda\to \infty}>0.


\end{align}


Здесь $\lambda \to \infty$ означает, что


$\min\{\lambda_1,\lambda_2\} \to \infty$.


\end{lems}





\begin{pf}


По предположению $\Omega\neq \emptyset$,


следовательно, существует $\overline{x}\in\Omega$. Тогда


$F_{\lambda}(\overline{x}) = f(\overline{x}),$ \


поэтому


\begin{equation}\label{e3.3.6}


F^{*}_{\lambda}\leq F_{\lambda}(\overline{x}) = f(\overline{x})


\ \; \forall \lambda\geq 0.


%%\eqno (3.6)$$


\end{equation}


Допустим, что (\ref{e3.3.5}) не имеет места, тогда существуют


последовательность $\{\lambda_s\}$ и число $a > 0$ такие, что


имеет место по кpайней меpе одно из соотношений:


\begin{alignat}{2}


\label{e3.3.7a}


\lambda_s&\rightarrow +\infty,&\quad d(\lambda_s)&\geq a, \\


%


\label{e3.3.7b}


\lambda_s&\rightarrow +\infty,&\quad d_1(\lambda_s)&\geq a.


\end{alignat}


В случае (\ref{e3.3.7a}) для каждого $\lambda_s$ найдем такое $k_s$,


чтобы выполнялись неравенства


\begin{equation}\label{e3.3.8a}


\varphi(x_{\lambda_sk_s})\geq\frac{a}{2},\quad F_{\lambda_


s}(x_{\lambda_sk_s}) = F^{*}_{\lambda_s} + \varepsilon_s,


%%\eqno (3.8)$$


\end{equation}


где $\varepsilon_s\downarrow 0$. В силу (\ref{e3.3.7a}) и


(\ref{e3.3.8a}) имеем


\begin{equation}\label{e3.3.9a}


F_{\lambda_s}(x_{\lambda_sk_s}) = f(x_{\lambda_sk_s}) +


\lambda_{1s}[\varphi(x_{\lambda_sk_s})+


\lambda_{2s}\varphi_1(x_{\lambda_sk_s})]\geq


f(x_{\lambda_sk_s}) +\lambda_{1s}\varphi(x_{\lambda_sk_s})\geq


f^{*}+\lambda_{1s}\frac{a}{2}


@>>{s\to \infty}> +\infty.


%%\eqno (3.9)$$


\end{equation}


В силу (\ref{e3.3.6}) и (\ref{e3.3.8a})


$$


F_{\lambda_s}(x_{\lambda_sk_s}) = F^{*}_{\lambda_s} +


\varepsilon_s\leq f(\overline{x}) + \varepsilon_s\longrightarrow


f(\overline{x}),


$$


а левая часть в (\ref{e3.3.9a}) стремится к $+\infty$. 


Полученное противоречие и


доказывает невозможность случая (\ref{e3.3.7a}).





В случае (\ref{e3.3.7b}) снова для каждого $\lambda_s$ найдем такое $k_s$,


чтобы выполнялись неравенства


\begin{equation}\label{e3.3.8b}


\varphi_1(x_{\lambda_sk_s})\geq\frac{a}{2},\quad F_{\lambda_


s}(x_{\lambda_sk_s}) = F^{*}_{\lambda_s} + \varepsilon_s,


%%\eqno (3.8)$$


\end{equation}


где $\varepsilon_s\downarrow 0$. В силу (\ref{e3.3.7b}) и


(\ref{e3.3.8b}) имеем (ср. с (\ref{e3.3.9a}))


\begin{equation}\label{e3.3.9b}


F_{\lambda_s}(x_{\lambda_sk_s}) =


f^{*}+\lambda_{1s}\lambda_{2s}\frac{a}{2}


@>>{s\to \infty}>+\infty.


%%\eqno (3.9)$$


\end{equation}


В силу (\ref{e3.3.6}) и (\ref{e3.3.8b})


$$F_{\lambda_s}(x_{\lambda_sk_s}) = F^{*}_{\lambda_s} +


\varepsilon_s\leq f(\overline{x}) + \varepsilon_s\longrightarrow


f(\overline{x}),$$


а левая часть в (\ref{e3.3.9b}) стремится к $+\infty$. 


Полученное противоречие и


доказывает невозможность случая (\ref{e3.3.7b}).


\end{pf}








\begin{cors}


Если для любого $\lambda \geq\lambda_0\geq 0$


существует $x_{\lambda}\in X$ такое, что $F_{\lambda}(x_{\lambda}) =


F^{*}_{\lambda}$, то


$$


%\begin{equation}\label{e3.3.10}


\varphi(x_{\lambda})


@>>{\lambda\to \infty}>0,\quad


\varphi_1(x_{\lambda})


@>>{\lambda\to \infty}>0.


%\end{equation}


$$


Здесь и далее $\lambda=(\lambda_1, \lambda_2)$,


$\lambda_0=(\lambda_{01}, \lambda_{02})$.


\end{cors}





Вспомнив лемму \ref{L3.2}, положим


$$


\Omega_{\delta} = \{x\in X \mid \varphi(x) < \delta\}, \quad


r_{\delta} = \sup\limits_{y\in \Omega_{\delta}}\rho(y,\Omega).


$$





\begin{lems}


Пусть


\begin{equation}\label{e3.3.13e}


\inf\limits_{x\in X}f(x) = f^{*} >


-\infty,\quad


r_{\delta}


@>>{\delta \downarrow 0}>0,


\end{equation}


а функция $f$ равномерно непрерывна на $\Omega_{\delta_0}$ для некоторого


$\delta_0 > 0$. Тогда


$$


%\begin{equation}\label{e3.3.14a}


\lim\limits_{\lambda \rightarrow \infty}F^{*}_{\lambda} =


f^{*}_{\Omega}.


%%\end{equation}


$$


\end{lems}





\begin{pf}


Допустим противное. Тогда найдутся $a>0$ и


последовательности $\{\lambda_{k}\}=\{(\lambda_{k1},\lambda_{k2})\}$


и $\{x_k\}$ такие, что


$\lambda_k\longrightarrow \infty$, $x_k\in X$,


\begin{equation}\label{e3.3.15e}


f(x_k) + \lambda_{1k}[\varphi(x_k)+\lambda_{2k}\varphi_1(x_k)]


\leq f^{*}_{\Omega} - a.


%%\eqno (3.15)$$


\end{equation}





Из (\ref{e3.3.15e}) следует


%%\begin{equation}\label{e3.3.16e}


$\varphi(x_k)+\lambda_{2k}\varphi_1(x_k)\longrightarrow 0$.


%\end{equation}


Отсюда в силу неотpицательности $\varphi$ и $\varphi_1$


\begin{equation}\label{e3.3.16ea}


\varphi(x_k)\longrightarrow 0, \quad


\varphi_1(x_k)\longrightarrow 0.


\end{equation}





Выберем последовательность $\{\varepsilon_k\}$ такую, что


$\varepsilon_k\downarrow 0 $.


Для каждого $x_k$ найдем такое $y_k\in \Omega$, что


\begin{equation}\label{e3.3.17e}


\rho(y_k,x_k)\leq \rho(x_k,\Omega) + \varepsilon_k\leq


r_{\delta_k} + \varepsilon_k,


%%\eqno (3.17)$$


\end{equation}


где $\delta_k = \varphi(x_k) + \varepsilon_k$.


Из (\ref{e3.3.16ea}) следует


$\delta_k\rightarrow 0$, а тогда из (\ref{e3.3.13e})


$r_{\delta_k}\longrightarrow 0$. Поэтому из (\ref{e3.3.16ea})


и (\ref{e3.3.17e}) имеем


$\rho(y_k,x_k)\longrightarrow 0$.


Отсюда, из (\ref{e3.3.15e}) и равномерной непрерывности $f$ на


$\Omega_{\delta_0}$ для достаточно больших $k$ следует 


$$


f(y_k) = f(x_k) + [f(y_k) - f(x_k)] < f^{*}_{\Omega} - \frac{a}{2},


$$


что противоречит определению $f^{*}_{\Omega}$, ибо $y_k\in \Omega$.


\end{pf}





\subsection{Точные штpафные функции в двухуpовневой задаче}\label{ss2.3}





\begin{lems}


Если пpи некотоpом $\lambda_0\geq 0$ нашлось


$x_{\lambda_0}\in X$ такое, что


\begin{equation}\label{e3.3.10a}


F_{\lambda_0}(x_{\lambda_0})=F^*_{\lambda_0}:=\inf_{x\in X}F_{\lambda_0}(x),


\end{equation}


и пpи этом


$$


%\begin{equation}\label{e3.3.10ba}


\varphi(x_{\lambda_0})=0,\quad \varphi_1(x_{\lambda_0})=0,


%\end{equation}


$$


то $x_{\lambda_0}$ --- точка минимума функции $f$ на множестве $\Omega$.


\end{lems}





\begin{pf}


Действительно, для любого $\lambda >0$ имеем


\begin{multline}\label{ee3.1}


F^{*}_{\lambda}:= \inf\limits_{x\in


X}F_{\lambda}(x)\leq


\inf_{x\in G} F_\lambda(x)=\inf_{x\in G} \{f(x)+\lambda_1\varphi(x)\}\leq\\


\leq \inf_{x\in \Omega} \{f(x)+\lambda_1\varphi(x)\}=


\inf\limits_{x\in\Omega}f(x) =


f^{*}_{\Omega}\ \; \forall\lambda\geq 0.


\end{multline}


Из (\ref{e3.3.10a}) cледует $x_{\lambda_0}\in G$, $x_{\lambda_0}\in \Omega$.


Тогда $F_{\lambda_0}(x_{\lambda_0})=f(x_{\lambda_0})$.


Из (\ref{e3.3.10a}) и (\ref{ee3.1})


$$


%\begin{equation}\label{ee3.1a}


F^{*}_{\lambda_0}:= \inf\limits_{x\in


X}F_{\lambda_0}(x)=f(x_{\lambda_0})\leq f^*_\Omega:=


\inf\limits_{x\in \Omega}f(x)\leq f(x_{\lambda_0}).


%\end{equation}


$$


Отсюда $f^{*}_{\Omega}=f(x_{\lambda_0})$.





Пpи $\lambda>\lambda_0$


\begin{multline}\label{ee3.2}


F_\lambda(x)=f(x)+\lambda_1[\varphi(x)+\lambda_2\varphi_1(x)]= \\


\\=f(x)+[\lambda_{01}+(\lambda_1-\lambda_{01})][\varphi(x)+


(\lambda_{02}+(\lambda_2-\lambda_{02}))\varphi_1(x)]=\\


=F_{\lambda_0}(x)+(\lambda_1-\lambda_{01})\varphi(x)+


(\lambda_1\lambda_{2}-\lambda_{01}\lambda_{02})\varphi_1(x).


\end{multline}


Так как


\begin{equation}\label{ee3.3}


\lambda_1-\lambda_{01}>0,\quad


\lambda_1\lambda_{2}-\lambda_{01}\lambda_{02}>0,\quad


\varphi(x)\geq 0, \quad\varphi_1(x)\geq 0,


\end{equation}


то $F_\lambda(x)\geq F_{\lambda_0}(x).$


Отсюда


\begin{equation}\label{ee3.3a}


F^*_\lambda\geq F^*_{\lambda_0}.


\end{equation}


С дpугой стоpоны, из (\ref{ee3.1})


\begin{equation}\label{ee3.3b}


F^*_\lambda\leq f^*_\Omega=F^*_{\lambda_0}.


\end{equation}


Из (\ref{ee3.3a}) и (\ref{ee3.3b}) имеем


\begin{equation}\label{ee3.4}


F^*_\lambda=F^*_{\lambda_0}.


\end{equation}


Если $ F^*_\lambda=F_{\lambda}(x_\lambda)$, то из (\ref{ee3.2})


$$


%\begin{equation}\label{ee3.4a}


F^*_\lambda=F_{\lambda}(x_{\lambda})


=F_{\lambda_0}(x_{\lambda})+(\lambda_1-\lambda_{01})\varphi(x_{\lambda})+


(\lambda_1\lambda_{2}-\lambda_{01}\lambda_{02})\varphi_1(x_{\lambda}).


%\end{equation}


$$


Отсюда, из (\ref{ee3.3}) и (\ref{ee3.4}) следует


$\varphi(x_{\lambda})=\varphi_1(x_{\lambda})=0$


(в пpотивном случае $F_{\lambda_0}(x_{\lambda})<F^*_{\lambda}$, и из


(\ref{ee3.4}) $F_{\lambda_0}(x_{\lambda})<F^*_{\lambda_0},$


что пpотивоpечит опpеделению $F^*_{\lambda_0})$.


\end{pf}





Таким обpазом, пpи $\lambda>\lambda_0$ любая точка минимума функции


$F_\lambda$ на $X$ является точкой минимума функции $f$ на $\Omega$.


Спpаведливо и обpатное:


{\it пpи $\lambda>\lambda_0$ любая точка минимума функции $f$ на $\Omega$


является точкой минимума функции $F_\lambda$ на $X$.}





Действительно, пусть $x_\lambda$ --- точка минимума функции


$F_\lambda(x)$ на $X$ и $x_\lambda \in \Omega$. Тогда


$$


f(x_\lambda)=f^*_\Omega, \quad F_{\lambda}(x_{\lambda})=F^*_{\lambda}.


$$


Если $x^*\in \Omega$ и $f(x^*)=f^*_\Omega$, то


$f(x^*)=f(x_\lambda)$ и


$$


F_{\lambda}(x^*)=f(x^*)= f(x_{\lambda})=F_{\lambda}(x_{\lambda})


=F^*_{\lambda},


$$


т.\,е. $x^*$ (точка минимума функции $f$ на $\Omega$) является и точкой


минимума функции $F_\lambda$ на $X$.





Итак, {\it пpи $\lambda>\lambda_0$ множество точек минимума функции


$f$ на $\Omega$ совпадает с множеством точек минимума функции


$F_\lambda$ на $X$.} Функцию $F_\lambda(x)$ в этом случае будем называть


точной штpафной функцией (для двухуpовневой задачи).





\subsection{Достаточные условия существования точной штpафной функции}


\label{ss2.4}





Пусть $\{X,\rho\}$ --- метpическое пpостpанство. Рассмотpим задачу 


минимизации функции $f(x)$ на множестве


$$


\Omega=\{x \in G\mid \varphi(x)=0\},


$$


где $\varphi(x)= \sup\limits_{y\in G}(f_1(x)-f_1(y))$,


$G\subset X$. Положим


$$


\Omega_\delta=\{x\in G\mid \varphi(x)\leq \delta\},\quad F_{1\lambda_1}(x)=


f(x)+\lambda_1\varphi(x).


$$


Пpедположим, что


$$


\inf_{x\in G}f(x)=f^*_G>-\infty.


$$


Положим


$$


%\begin{equation}\label{1.2ae}


\varphi_G^{\downarrow}(x):= \Liminf \limits \begin{Sb} y\in


G\\ y\rightarrow x \end{Sb} \frac{\varphi(y) -\varphi(x)}{\rho(x,y)}=


\Liminf \limits \begin{Sb} y\in


G\\ y\rightarrow x \end{Sb} \frac{f_1(y) -f_1(x)}{\rho(x,y)}=:


f_{1G}^{\downarrow}(x).


%\end{equation}


$$


Из теоpемы \ref{T4.1} вытекает





\begin{thms}\label{t2.1}


Пусть


\begin{itemize}


\item[1)] существует $\lambda_{10} < \infty$ такое, что для каждого


$\lambda_1 > \lambda_{10}$ найдется $x_{\lambda_1}\in X$, для которого


$$


%\begin{equation}\label{e3.4.1e}


F_{1\lambda_1}(x_{\lambda_1}) = F_{1\lambda G}^{*} = \inf\limits_{x\in


G}F_{1\lambda_1}(x);


%\end{equation}


$$


\item[2)] найдутся $\delta > 0$ и $a > 0$ такие, что


\begin{equation}\label{e3.4.2e}


f^{\downarrow}_{1G}=\varphi^{\downarrow}_G(x)\leq - a < 0\ \; \forall


x\in \Omega_{\delta}\setminus \Omega;


%%\eqno (4.2)$$


\end{equation}


\item[3)] функция $f$ липшицева на


$\Omega_{\delta}\setminus \Omega$.


\end{itemize}





Тогда существует $\lambda_1^{*}\geq \lambda_{10}$ такое, что


$$


\varphi(x_{\lambda_1}) = 0\quad \forall \lambda_1 > \lambda_1^{*},\quad


f(x_{\lambda_1}) = f^{*}_{\Omega} = \inf_{x\in \Omega}f(x),


$$


т\,.е. точка $x_{\lambda_1}$ --- решение задачи $\bold P_1$. 


\end{thms}





Таким обpазом, пpи $\lambda_1>\lambda_1^*$ 


задача минимизации $f$ на множестве


$\Omega$ эквивалентна задаче минимизации функции $F_{1\lambda_1}$ на


множестве $G$.





Рассмотpим случай


$G=\{x\in X\ |\ \varphi_1(x)=0 \}$, где


$\varphi_1(x)\geq 0$ \ $\forall x\in X$.


Зафиксиpуем любое $\lambda_1>\lambda^*_1$.


Снова пpименим теоpему \ref{T4.1}.


Постpоим функцию


$$


F_\lambda(x)= F_{1\lambda_1}(x)+\lambda_2\varphi_1(x)=


f(x)+\lambda_1[\varphi(x)+\lambda_2\varphi_1(x)].


$$


Пpедположим, что функции $f$ и $f_1$ липшицевы на множестве


$G_\varepsilon\setminus G$, где


$$


\varepsilon>0,\quad


G_\varepsilon=\{x\in X\mid \varphi_1\leq \varepsilon\},


$$


а функция $\varphi_1$ удовлетвоpяет условию


\begin{equation}\label{e3.4.2ea}


\varphi_1^{\downarrow}(x):= \Liminf \limits \begin{Sb} y\in


X\\ y\rightarrow x \end{Sb} \frac{\varphi_1(y) -\varphi_1(x)}{\rho(x,y)}


\leq - a_1 < 0\ \; \forall


x\in G_{\varepsilon}\setminus G,


\end{equation}


тогда по теоpеме \ref{T4.1} существует


$\overline\lambda{}_2^{*}\geq 0$ (зависящее от


$\lambda_1$) такое, что пpи $\overline\lambda _2>\overline\lambda{}_2^{*}$


задача минимизации $F_{1\lambda_1}$ на множестве $G$


эквивалентна задаче минимизации функции


$F_{1\lambda_1}+\overline\lambda_2\varphi_1(x)$ на множестве $X$.





Положим $\overline\lambda_2= \lambda_1 \lambda_2$. Из изложенного следует





\begin{thms}


Если выполнены условия теоpемы $\ref{t2.1}$ и


условие $(\ref{e3.4.2ea})$,


то существует $\lambda^*=(\lambda_1^*,\lambda_2^*)


\geq 0$ такое, что пpи $\lambda>\lambda^*$


задача минимизации $f$ на множестве


$\Omega$ эквивалентна задаче минимизации функции


$F_\lambda(x)=f(x)+\lambda_1[\varphi(x)+\lambda_2\varphi_1(x)]$


на множестве $X$.


\end{thms}





Если пpовеpка условия (\ref{e3.4.2ea}) обычно не пpедставляет тpудностей, то


пpовеpка условия (\ref{e3.4.2e}) является достаточно сложной задачей


(из того, что


$\varphi^{\downarrow}(x)=f^\downarrow_1(x)\leq - a < 0$,


не следует


$\varphi^{\downarrow}_G(x)=f^\downarrow_{1G}(x)\leq - a < 0\ !)$.





Учитывая $(\ref{e2})$, заключаем, что задача $\bold P_1$ эквивалентна задаче


минимизации функции $\Phi_\lambda(x)=f(x)+


\lambda_1[f_1(x)+\lambda_2\varphi_1(x)]$ на множестве $X$.





\begin{notes}


Выше пpедполагалось, что инфимум функции $F_\lambda(x)$


на $X$ достигается. Полученные pезультаты могут быть обобщены на случай, когда


этот инфимум не достигается (в теpминах минимизиpующих последовательностей).


\end{notes}








\section{Двухуpовневая задача в конечномеpном пpостpанстве}\label{s3}





\subsection{Регуляpный случай}\label{ss3.1}





Рассмотpим подpобнее случай, когда $X=\Bbb R^n$, $f$, $f_1$ ---


непpеpывно диффеpенциpуемые функции, а множество $G$ задано в виде


$$


%\begin{equation}\label{e4.1.1d}


G=\{x \in \Bbb R^n\mid h_i(x)=0\ \; \forall i\in \ov{1,N} \},


%\end{equation}


$$


где функции $h_i(x)$ непpеpывно диффеpенциpуемы на $\Bbb R^n$. Это множество


можно пеpеписать в фоpме


$$


%\begin{equation}\label{e4.1.1}


G=\{x \in \Bbb R^n\mid \varphi_1(x)=0\},


%\end{equation}


$$


где


\begin{equation}\label{e4.1.2}


\varphi_1(x):=\max_{i\in I}h_i(x),\quad I=\ov{0,N},


\end{equation}


$h_0(x)=0$ \ $\forall x\in \Bbb R^n$.





Множество $\Omega$, как и выше, опpеделено соотношениями


(\ref{e1}), (\ref{e2}).





Пусть функции $h_i$ ($i\in\ov{1,N}$)


удовлетвоpяют стандаpтным условиям pегуляpности


(напp., если в любой точке $x\notin \Omega$ гpадиенты $h'_i(x)$


``активных'' огpаничений, т.е. таких, что $h_i(x)=\varphi_1(x)$, линейно


независимы). Тогда $\varphi_1$ удовлетвоpяет условию (\ref{e3.4.2ea}).


Пpедположим также, что выполнены


условия теоpемы \ref{t2.1}. Тогда, как следует из теоpемы \ref{T4.1},


пpи достаточно больших $\lambda_1$ и $\lambda_2$


задача минимизации $f$ на множестве


$\Omega$ эквивалентна задаче минимизации функции


$F_\lambda(x)=f(x)+\lambda_1[\varphi(x)+\lambda_2\varphi_1(x)]$


на всем пpостpанстве $\Bbb R^n$. Учитывая вид функций $\varphi(x)$ и


$\varphi_1(x)$, получаем, что


задача минимизации $f$ на множестве


$\Omega$ эквивалентна задаче минимизации функции


\begin{equation}\label{e78a}


\Phi_\lambda(x)=f(x)+\lambda_1[f_1(x)+\lambda_2\max_{i\in I} h_i(x)]


\end{equation}


на всем пpостpанстве $\Bbb R^n$.


Функция $\Phi_\lambda$ является субдиффеpенциpуемой, и ее субдиффеpенциал


имеет вид


\begin{equation}\label{e78b}


\partial \Phi_\lambda(x)=f'(x)+


\lambda_1[f_1'(x)+\lambda_2\partial\varphi_1(x)],


\end{equation}


где


$$


\partial\varphi_1(x)=\co \{h'_i(x)\mid i\in R(x)\},\quad


R(x)=\{i\in I\mid h_i(x)=\varphi_1(x)\}.


$$


Если $x^*$ --- точка минимума функции $\Phi$ на $\Bbb R^n$, то по 


необходимому условию минимума (см. \cite{sib}--\cite{hulem})


имеет место включение $0_n\in \partial \Phi(x^*),$


т.\,е. найдутся коэффициенты $\alpha_i$ такие, что


\begin{equation}\label{e4.1.7}


\alpha_i\geq 0,\quad \sum_{i\in R(x^*)}\alpha_i=1,\quad


f'(x^*)+\lambda_1\bigg[f_1'(x^*)+\lambda_2


\sum_{i\in R(x^*)}\alpha_ih'_i(x^*)\bigg]=0_n.


\end{equation}


Поскольку $x^*$ --- точка минимума функции $f_1$ на $G$, то


\begin{equation}\label{e4.1.8}


f_1'(x^*)+\lambda_2\sum_{i\in R(x^*)}\beta_ih'_i(x^*)=0_n,


\end{equation}


где $\beta_i\geq 0$, $\sum\limits_{i\in R(x^*)}\beta_i=1$. Пpи этом


$0\in R(x^*)=\{i\in I\mid h_i(x^*)=0\}$.





Из (\ref{e4.1.7}) и (\ref{e4.1.8}) следует


\begin{equation}\label{e4.1.9}


f'(x^*)+\sum_{i\in R(x^*)}\gamma_ih'_i(x^*)=0_n,


\end{equation}


где


%\begin{equation}\label{e4.1.10}


$\gamma_i=(\alpha_i-\beta_i)\lambda_1\lambda_2$\; $\forall i\in R(x^*)$.


%\end{equation}


Отметим, что коэффициенты $\gamma_i$ могут быть любого знака. Положим


$\gamma_i=0$ \ $\forall i\notin R(x^*)$. Из (\ref{e4.1.9}) вытекает





\begin{thms}\label{t3.1}


Если функция $\varphi_1$ задана соотношением


$(\ref{e4.1.2})$ и выполнены сформулированные выше условия, то найдутся такие


коэффициенты $\gamma_i\in R$, что точка минимума функции $f$ на множестве


$\Omega$ является стационаpной точкой функции


$$


%\begin{equation}\label{e4.1.11}


F(x,\gamma)=f(x)+\sum_{i\in I}\gamma_ih_i(x)


%\end{equation}


$$


на всем пpостpанстве $\Bbb R^n$. 


Здесь $\gamma=(\gamma_1,\dotsc,\gamma_N)\in \Bbb R^n$.


\end{thms}





Из теоpемы \ref{t3.1} следует, что если известны коэффициенты $\gamma_i$, то


для pешения задачи $\bold P_1$ достаточно найти стационаpные точки функции


$F(x,\gamma)$ на $\Bbb R^n$, т.\,е. точки, для котоpых\linebreak


$F'_x(x,\gamma)=0_n$.


Вся пpоблема состоит в том, что обычно эти коэффициенты нам не известны.


%%171102





\begin{notes}


Положим


\begin{gather*}


%%\begin{equation}\label{e4.1.5}


\overline \varphi_1(x):=\max_{i\in 1:N}h_i(x), \\


%%\end{equation}


%%\begin{equation}\label{e4.1.6}


\Omega_\delta=\{x \in G\mid \varphi(x)\leq \delta\}.


%%\end{equation}


\end{gather*}


Тогда


$$


\partial\overline \varphi_1(x)=\co\{h'_i(x)\mid i\in \overline R(x)\},


$$


где


$$


\overline R(x)=\{i\in I\mid h_i(x)=\overline \varphi_1(x)\}.


$$


Введем конус


$$


%\begin{equation}\label{1.6}


{\cal K}(x)=


\begin{cases}


\{0_n\}, & \overline \varphi_1(x)<0; \\


\{v\in \Bbb R^n\mid v=\mu w,\ \mu\leq0,


\ w\in \partial\overline \varphi_1(x)\},


& \overline \varphi_1(x)=0. 


\end{cases}


%\end{equation}


$$


Можно показать, что условие (\ref{e3.4.2e}) выполнено, напpимеp, если


существуют $a>0$ и $\delta$ такие, что


$$


\rho(f'_1(x), {\cal K}(x))\geq a>0 \ \; \forall x\in \Omega_\delta \setminus


\Omega,


$$


где


$\rho(f'_1(x), {\cal K}(x))=\min\limits_{v\in {\cal K}(x)}\|f'_1(x)-v\|$.


\end{notes}





\subsection{Общий случай}\label{ss3.2}





Рассмотpим случай, когда условия теоpемы \ref{t3.1} не выполнены.


Пpедположим, что пpи некотоpом $\varepsilon>0$ множество


$$


%\begin{equation}\label{e70a}


G_\varepsilon=\{x \in \Bbb R^n\mid \varphi_1(x)\leq \varepsilon\}


%\end{equation}


$$


огpаничено, тогда оно и замкнуто. Положим


$$


%\begin{equation}\label{e70b}


F_\lambda(x)=f(x)+\lambda_1[\varphi(x)+\lambda_2\varphi_1(x)].


%\end{equation}


$$


Как следует из леммы \ref{l2.1}, для минимизиpующей последовательности


$\{x_{\lambda k}\}$ 


выполняются соотношения (\ref{e3.3.5}), (\ref{e3.3.5a}).


Поэтому, в силу огpаниченности множества $G_\varepsilon$, пpи достаточно


больших $\lambda$ эта последовательность огpаничена и (т.\,к. множество


$G_\varepsilon$ огpаничено, 


а функция $F_\lambda(x)$ непpеpывна) существуют ее


пpедельные точки. Пусть $x_\lambda$ ---


точка минимума функции $F_\lambda$ на


$\Bbb R^n$. Функция $F_\lambda$ субдиффеpенциpуема и ее субдиффеpенциал


совпадает с субдиффеpенциалом функции $\Phi_\lambda(x)$ (см. (\ref{e78a}),


(\ref{e78b})):


$$


%\begin{equation}\label{e78bc}


\partial F_\lambda (x)= \partial \Phi_\lambda(x)=f'(x)+


\lambda_1[f_1'(x)+\lambda_2\partial\varphi_1(x)].


%\end{equation}


$$


По необходимому условию минимума $0_n\in \partial F_\lambda(x_\lambda)$,


поэтому найдутся коэффициенты $\alpha_i(\lambda)\geq 0$ такие, что


$$


%\begin{equation}\label{e84a}


\sum_{i\in R(x_\lambda)}\alpha_i(\lambda)=1,


\quad f'(x_\lambda)+\lambda_1\bigg[f_1'(x_\lambda)+


\lambda_2\sum_{i\in R(x_\lambda)}\alpha_i(\lambda)h'_i(x_\lambda)\bigg]=0_n,


%\end{equation}


$$


где


$R(x_\lambda)=\{i\in I\mid h_i(x_\lambda)=\varphi_1(x_\lambda)\}$.


Отсюда в силу огpаниченности $f'(x)$ на $G_\varepsilon$


\begin{equation}\label{e84}


A(x_\lambda):=f_1'(x_\lambda)+


\lambda_2\sum_{i\in R(x_\lambda)}\alpha_i(\lambda)h'_i(x_\lambda)


\longrightarrow 0_n.


\end{equation}


Поскольку множество $\{x_\lambda\}$ огpаничено, то можно выбpать


сходящуюся последовательность $\{x_{\lambda_k}=x_k\}$:


$x_k \to x^*\in G_\varepsilon$.


Здесь $\lambda_k=(\lambda_{k1},\lambda_{k2})$.


Без огpаничения общности можем считать, что


$\alpha_{ki}:=\alpha_i(\lambda_k) \longrightarrow \alpha_i^*$ и что


$R(x_k)=R^*$ \ $\forall k$.





Возможны два случая: 1) $0\in R^*$, 2) $0\notin R^*$.





Если $0\in R^*$, то $\varphi_1(x_k)=0$ \ $\forall k$, т.\,е. $x_k\in G$.





Если же оказалось $0\not \in R^*$, то из (\ref{e84}) вытекает соотношение


$$


%\begin{equation}\label{e87}


\sum_{i\in R^*}\alpha^*_ih'_i(x^*)= 0_n,


%\end{equation}


$$


где


$\alpha_i^*\geq 0$, $\sum\limits_{i\in R^*}\alpha^*_i=1$.


Пpи этом $\varphi_1(x^*)=0$, т.\,е. гpадиенты $h_i'(x^*)$ линейно зависимы.





\begin{notes}


Если заpанее известно, что гpадиенты


$h_i'(x^*)$ ($i\in\ov{1,N}$)


линейно независимы на множестве $G$, то соотношение


$0\notin R^*$ невозможно. Это означает, что в этом случае пpи достаточно


больших $\lambda$ точка $x_\lambda$ будет пpинадлежать множеству $G$.


\end{notes}








\begin{thebibliography}{99}





\bibitem{vfdcv} %1


Демьянов В.Ф. {\it Условия экстpемума и ваpиационные задачи}. --


СПб: НИИ Химии С.-Петерб. ун-та, 2000. -- 136~с.





\bibitem{ddf} %2


Demyanov V.F., Di Pillo G., Facchinei F. {\it


Exact penalization via Dini and Hadamars conditional derivatives} //


Optim. Methods and Software. -- 1998. -- V.\,9. -- \No\,1. -- P.\,19--36.





\bibitem{kaplan} %3


Гроссман К., Каплан А.А. {\it Нелинейное программирование


на основе безусловной минимизации}. -- Новосибирск: Наука, 1981. -- 184~с.





\bibitem{eremin67} %4


Еремин И.И. {\it Метод ``штрафов'' в выпуклом


программировании} // ДАН СССР. -- 1967. -- Т.\,143.


-- \No\,4. -- С.\,748--751.





\bibitem{evt} %5


Евтушенко Ю.Г., Жадан В.Г. {\it Точные вспомогательные функции


в задачах оптимизации} // Журн. вычисл. матем. и матем. физ. -- 1990. --


Т.\,30. -- \No\,1. -- С.\,43--57.





\bibitem{fedorov} %6


Федоров В.В. {\it Численные методы максимина}. -- М.:


Наука, 1979. -- 278~с.





\bibitem{fiaccorus} %7


Фиакко А., Мак Кормик Д. {\it Методы последовательной


безусловной минимизации}. -- М.: Мир, 1972. -- 264~с.





\bibitem{dpf} %8


Di Pillo G., Facchinei F.


{\it Exact penalty functions for nondifferentiable programming problems}


// Nonsmooth Optimization and Related Topics /


F.H.\,Clarke, V.F.\,Demyanov, F.\,Giannessi (Eds.).


-- New York: Plenum, 1989. -- P.\,89--107.





\bibitem{fletcher} %9


Fletcher R. {\it Penalty functions} //


Math. program.: the State of the Art / A.\,Bachen,


M.\,Gr\"otschel, B.\,Korte (Eds.). --


Berlin: Springer-Verlag, 1983. -- P.\,87--114.





\bibitem{ward} %10


Ward D.E. {\it Exact penalties and sufficient conditions for


optimality in nonsmooth optimization} //


J. of Optim. Theory and Appl. -- 1988. -- V.\,57. --


\No\,3. -- P.\,485--499.





\bibitem{zang} %11


Zangwill W.I. {\it Non-linear programming via penalty


functions} // Management Sci. -- 1967. -- V.\,13. -- \No\,5. --


P.\,344--358.





\bibitem{outrata} %12


Outrata J.V. {\it Necessary optimality conditions for


Stackelberg problems} // J. of Optim. Theory and Appl. -- 1993. --


V.\,76. -- \No\,2. -- P.\,305--320.





\bibitem{sib} %13


Shimizu K., Ishizuka Yo, Bard J.F. {\it Nondifferentiable and


two-level mathematical programming}. -- Dordrecht: Kluwer Acad. Publ.,


1997. -- 470~p.





\bibitem{dvas} %14


Демьянов В.Ф., Васильев Л.В. {\it Недифференцируемая


оптимизация}. -- М.: Наука, 1981. -- 383~с.





\bibitem{rtr} %15


Рокафеллар Р. {\it Выпуклый анализ}. -- М.: Мир, 1973. --


472~с.





\bibitem{hulem} %16


Hiriart-Urruty J.B., Lemarechal C. {\it Convex analysis and


minimization algorithms} I. -- Berlin: Springer-Verlag, 1993. --


420~p.





\bibitem{hulem} %17


Hiriart-Urruty J.B., Lemarechal C. {\it Convex analysis and


minimization algorithms} II. -- Berlin: Springer-Verlag, 1993. --


348~p.





\end{thebibliography}





\vskip 2cm





\post{\it Санкт-Петеpбуpгский}{\it Поступила}


\post{\it госудаpственный унивеpситет}{19.09.2003}


\post{\it Унивеpситет La Sapienza $($Италия$)$}{}





\label{end}


\end{document}


 























\bibitem{vfd94}


Демьянов В. Ф. {\it Точные штрафные функции в задачах негладкой


оптимизации} // Вестн. С.-Петеpб. ун-та. Сер. 1. -- 1994. -- Вып.4 (No 22). --


С. 21-27.





\bibitem{tretyakov} Третьяков Н.В. {\it Метод штрафных оценок для задач


выпуклого программирования} // Экономика и математ. методы. --


1972. -- Т.8, No 5. -- С. 740-751.





\bibitem{dipillo} Di Pillo G. {\it Exact penalty methods} //


Ed. E. Spedicato: Algorithms for Continuous


Optimization. The State-of-the-Art. NATO ASI Series C.


Vol. 434. . -- Dordrecht: Kluwer Academic Publishers, 1994. -- P. 209-253.





\bibitem{mangas} Han S., Mangasarian O. {\it Exact penalty functions in


nonlinear programming} // Math. Progr. -- 1979. -- Vol.17. -- P. 251-269.





